% !TeX root = ../document.tex
\section{Introduction}

The deployment of Artificial Intelligence (AI) on embedded systems, commonly referred to as Edge AI, represents a transformative frontier in modern computing. Unlike cloud-based AI, which relies on virtually unlimited computational resources, Edge AI requires the execution of sophisticated models on hardware with severe constraints in terms of memory, processing power, and energy consumption. Among the most prevalent platforms for these applications is the STM32 family of microcontrollers, which offers a versatile range of solutions for industrial, automotive, and consumer electronics.

However, the path from a high-level AI model developed in Python to an efficient firmware implementation in C remains fraught with complexity. Developers must navigate a fragmented ecosystem of tools—including STM32CubeMX for hardware configuration, STEdgeAI for model optimization, and various Python libraries for training—creating a significant "plumbing" problem. This disconnect between the AI development lifecycle and embedded systems engineering leads to slow iteration cycles, high error rates, and a steep barrier to entry for practitioners who lack deep expertise in both domains.

Current methodologies for Edge AI development often rely on manual, sequential steps that are difficult to reproduce and optimize. While conventional MLOps (Machine Learning Operations) has matured in cloud environments, providing automated pipelines for training and deployment, the embedded domain lacks comparable end-to-end infrastructure. Most existing automation tools focus on isolated tasks, such as quantization or code generation, leaving the broader orchestration of the development lifecycle—spanning firmware setup, model modification, and final hardware integration—as a manual and error-prone effort for the designer.

Recent breakthroughs in Large Language Models (LLMs) and agentic workflows offer a promising solution to this fragmentation. LLMs have demonstrated remarkable capabilities in code synthesis and reasoning across diverse programming domains. However, their potential has yet to be fully realized in the context of specialized embedded AI development, where architectural constraints and hardware-specific requirements demand more than simple code generation. There is a need for a system that can not only generate code but also orchestrate complex, multi-stage workflows and make informed decisions about model-hardware compatibility.

This thesis addresses this gap by proposing and implementing an end-to-end MLOps orchestration framework for STM32 AI projects. Rather than treating AI as an isolated component, the proposed system leverages a multi-agent architecture built on LangGraph to coordinate the entire development lifecycle. This framework acts as an intelligent assistant capable of managing firmware generation, performing specialized "Neural Network Surgery" to adapt models to hardware constraints, and automating the integration and verification process on the target device.

The proposed approach integrates the development lifecycle through several specialized branches. The system handles initial firmware project creation, provides automated AI analysis and quantization via the STEdgeAI suite, and implements a sophisticated customization branch for model refinement. A key innovation of this work is the integration of automated modification logic, allowing the system to replace unsupported layers or optimize architectures specifically for the target microcontroller's memory profile.

The research demonstrates that a structured, agent-based orchestration can bridge the accessibility gap between high-level AI design and low-level embedded implementation. By automating the "plumbing" between diverse tools and providing an interactive assistant that understands the hardware context, a foundation is established for more efficient, reproducible, and accessible Edge AI development practices.

The primary contributions of this work are threefold. First, a comprehensive multi-agent architecture for automated Edge AI development is presented, demonstrating how agentic orchestration can manage the end-to-end MLOps lifecycle on STM32 platforms. Second, a specialized "Neural Network Customization" framework is developed, enabling automated model surgery and hyperparameter optimization (via the NNI toolkit) to ensure compatibility with constrained hardware. Third, a dual-layer evaluation methodology is established: a qualitative assessment focusing on development efficiency and the reduction of cognitive load, and a quantitative validation using hardware-specific metrics (latency, memory footprint) and LLM-as-Judge frameworks (using \textit{DeepEval}) to ensure the reliability and faithfulness of the generated configurations.

